\documentclass[11pt]{article}
\usepackage[margin=1in]{geometry}
\usepackage{amsmath,amssymb,amsthm}
\usepackage{enumitem}
\usepackage{mathtools}
\usepackage{hyperref}
\usepackage{pgffor}

\newtheorem*{theorem*}{Theorem}

\title{MIT 6.1200J Problem Set 1 – Solutions}
\author{Nouman Rahman}
\date{29th August 2026}

\newcommand{\N}{\mathbb{N}}
\newcommand{\R}{\mathbb{R}}
\newcommand{\Z}{\mathbb{Z}}
\newcommand{\isOdd}{\mathrm{isOdd}}
\newcommand{\isPrime}{\mathrm{isPrime}}
\newcommand{\isEven}{\mathrm{isEven}}
\newcommand{\sixish}{\mathrm{sixish}}
\newcommand{\partlabel}[1]{%
    \noindent\makebox[2em][l]{\textbf{(#1)}}%
}

\begin{document}
\maketitle

\section*{Problem 1: Collaboration Policy}

\partlabel{a}
Collaboration is permitted in groups of up to 3 or 4, so a group of 10 is \textbf{too large} and must be divided into subgroups. Hence, in this scenario, the students are \textbf{violating the policy}.

\medskip

\partlabel{b}
Jointly developing ideas is allowed, but photographing/sharing the jointly created notes so they can be directly consulted during individual write-ups conflicts with the requirement that the write-up be completed individually \textbf{without directly consulting shared notes}. Hence, in this scenario, students are \textbf{violating the policy}.

\medskip

\partlabel{c}
Even though the students worked together, they developed \textit{a detailed solution} together, and one \textbf{dictated it to the others}. The policy allows jointly developing a broad outline, but the detailed proof/writeup must be independently produced in each student's own words and understanding. Therefore, in this scenario, students are \textbf{violating the course policy}.

\medskip

\partlabel{d}
The important issue is that one student has already solved the problem and then uses knowledge of the correct answer to identify the flaw in another student's work. The policy specifically \textbf{restricts offering an explanation} to someone who did not solve the problem with you. So, both students are \textbf{violating the course policy}.

\medskip

\partlabel{e}
The sounding-board/help aspect can be permissible, but allowing friends to skim the completed write-up is not: students should \textbf{not show their write-ups to peers}, and peers should not look at them. Therefore, the students are \textbf{violating the policy}.

\medskip

\partlabel{f}
Again, even though they collaborated on solving the problem, \textbf{exchanging completed write-ups} and examining them violates the rule against showing your write-up to peers and looking at peers' write-ups. Hence, this \textbf{violates the policy}.

\medskip

\partlabel{g}
Looking at a previous class's completed solution and copying it \textbf{is expressly prohibited}.

\medskip

\partlabel{h}
The policy says students should not look for/read completed solutions from outside sources, and copying from any source—\textbf{including books—is prohibited}.

\medskip

\partlabel{i}
Reading a completed solution in a book and then \textbf{explaining that solution to collaborators} is still using an outside completed solution as part of solving the problem. The collaboration rules do not make an outside solution permissible. Hence, this \textbf{violates the policy}.

\medskip

\partlabel{j}
Even though late students may read published Canvas solutions after they are released, they may \textbf{not consult those solutions while composing their own proofs}. Changing wording or sentence structure does \textbf{not make this permissible}.

\newpage
\section*{Problem 2: Set Formulas and Propositional Formulas}

\partlabel{a}
Verify that $(P \wedge \bar Q) \vee (P \wedge Q)$ is equivalent to $P$ using a truth table.

\begin{center}
\begin{tabular}{cc|c|c|c|c|c}
$P$ & $Q$ & $\bar Q$ & $P \wedge \bar Q$ & $P \wedge Q$ & $(P\wedge\bar Q)\vee(P\wedge Q)$ & $P$ \\
\hline
T & T & F & F & T & T & T \\
T & F & T & T & F & T & T \\
F & T & F & F & F & F & F \\
F & F & T & F & F & F & F \\
\end{tabular}
\end{center}

The column for $(P\wedge \bar Q)\vee(P\wedge Q)$ agrees with the column for $P$ in every row, so the two formulas are equivalent.

\bigskip

\partlabel{b}
Prove that $A = (A - B) \cup (A \cap B)$ for all sets $A, B$, by showing $x \in A$ iff $x \in (A-B)\cup(A\cap B)$ for all elements $x$.

\begin{proof}
Let $x$ be arbitrary.

($\Rightarrow$) Suppose $x \in A$. Either $x \in B$ or $x \notin B$.
\begin{itemize}
\item If $x \notin B$, then $x \in A$ and $x \notin B$, so $x \in A - B$, hence $x \in (A-B)\cup(A\cap B)$.
\item If $x \in B$, then $x \in A$ and $x \in B$, so $x \in A \cap B$, hence $x \in (A-B)\cup(A\cap B)$.
\end{itemize}
Either way, $x \in (A-B)\cup(A\cap B)$.

($\Leftarrow$) Suppose $x \in (A-B)\cup(A\cap B)$. Then $x \in A-B$ or $x \in A \cap B$.
\begin{itemize}
\item If $x \in A - B$, then by definition $x \in A$ (and $x \notin B$).
\item If $x \in A \cap B$, then $x \in A$ (and $x \in B$).
\end{itemize}
Either way, $x \in A$.

Since $x \in A$ iff $x \in (A-B)\cup(A\cap B)$ for arbitrary $x$, the two sets are equal: $A = (A-B)\cup(A\cap B)$.
\end{proof}

\emph{Relation to part (a).} This is the set-theoretic mirror of part (a). Substitute $P := $ ``$x \in A$'' and $Q := $ ``$x \in B$''. Then:
\begin{itemize}
\item $x \in A - B$ corresponds to $P \wedge \bar Q$ (in $A$, not in $B$),
\item $x \in A \cap B$ corresponds to $P \wedge Q$ (in $A$, in $B$),
\item $\cup$ corresponds to $\vee$.
\end{itemize}
So the set identity $A = (A-B)\cup(A\cap B)$ is exactly the membership-level statement $P \leftrightarrow \big((P\wedge \bar Q)\vee(P\wedge Q)\big)$ from part (a), and the case analysis in the proof above is just an unfolding of the truth table: the two cases $x\notin B$ / $x\in B$ are the two ways $Q$ can be false or true, matching the two rows of the truth table where $P$ is true.

\newpage
\section*{Problem 3: Predicate Practice}

\partlabel{a}
The predicate $\sixish(n)$, ``$n$ is a natural number whose first digit is $6$'':
\[
\sixish(n) :=\ \exists k \in \N.\ \big(6\cdot 10^{k} \le n < 7 \cdot 10^{k}\big).
\]

\partlabel{b}
``There are exactly two natural numbers $n$ such that $Q(n)$ is true'':
\[
\exists m, n \in \N.\ \Big( m \ne n \ \wedge\ \forall k \in \N.\big(Q(k) \leftrightarrow (k = m \vee k = n)\big) \Big).
\]

\partlabel{c}
Goldbach's Conjecture: ``every even integer greater than $2$ can be written as the sum of two primes.'' Following the required form $\forall n.\ W(n)$ with all the logic folded into $W$:
\[
\forall k \in \N.\ \Big[ \big(k > 2 \wedge \isEven(k)\big) \rightarrow \exists m, n \in \N.\big(\isPrime(m) \wedge \isPrime(n) \wedge k = m+n\big) \Big],
\]
where, since $\isEven$ was not predefined for us, we take $\isEven(k) := \neg \isOdd(k)$.


\partlabel{d}
``Every prime larger than $100$ is sixish'':
\[
\forall n \in \N.\ \Big[ \big(n > 100 \wedge \isPrime(n)\big) \rightarrow \sixish(n) \Big].
\]

\partlabel{e}
Is $R(n) := \exists b \in \N.\ (n = 2b+5)$ equivalent to $\isOdd(n)$?

\emph{Claim: No, $R$ is \underline{not} equivalent to $\isOdd$.}

\begin{proof}
Suppose, toward a contradiction, that $R(n) \leftrightarrow \isOdd(n)$ for all $n \in \N$. Take $n = 3$. Since $3$ is odd, $\isOdd(3)$ is true, so by our assumption $R(3)$ must also be true. But
\[
R(3) := \exists b \in \N.\ (3 = 2b+5),
\]
and solving $3 = 2b+5$ gives $b = -1 \notin \N$, so no such $b$ exists and $R(3)$ is false. This contradicts $R(3)$ being true. Hence our assumption was wrong, and $R$ is not equivalent to $\isOdd$.

(Intuitively, $R(n)$ only captures odd numbers $n \ge 5$, since $b$ ranges over $\N$; it misses the smaller odd numbers $1$ and $3$.)
\end{proof}

\section*{Problem 4: Induction and Contradiction}

\partlabel{a}
\emph{Setup.} $Q(x)$ is a predicate on $\R$ such that for all $x, y \in \R$, if $Q(x \cdot y)$ is true then $Q(x)$ or $Q(y)$ is true. We show the stronger property: for every integer $n \ge 1$ and all $x_1, \dots, x_n \in \R$, if $Q(x_1 x_2 \cdots x_n)$ is true then at least one of $Q(x_1), \dots, Q(x_n)$ is true.

\begin{proof}
Let $P(n)$ be the statement: for all $x_1, \dots, x_n \in \R$,
\[
Q(x_1 x_2 \cdots x_n) \rightarrow Q(x_1) \vee Q(x_2) \vee \cdots \vee Q(x_n).
\]
We prove $P(n)$ holds for all integers $n \ge 1$ by induction on $n$.

\textbf{Base case ($n=1$):} $P(1)$ says: for all $x_1 \in \R$, $Q(x_1) \rightarrow Q(x_1)$, which is trivially true.

\textbf{Inductive step:} Suppose $n \ge 1$ and $P(n)$ holds; we show $P(n+1)$ holds. Let $x_1, \dots, x_n, x_{n+1} \in \R$ be arbitrary, and suppose $Q(x_1 x_2 \cdots x_n x_{n+1})$ is true. Write
\[
y := x_1 x_2 \cdots x_n, \qquad z := x_{n+1},
\]
so that $x_1 \cdots x_n x_{n+1} = yz$. By hypothesis, $Q(yz)$ is true. Applying the given property of $Q$ (i.e., $P(1)$ is not what we need here --- we use the defining property of $Q$ directly, which is exactly $P(2)$ applied to $y,z$) to $y$ and $z$:
\[
Q(yz) \rightarrow Q(y) \vee Q(z).
\]
Since $Q(yz)$ holds, we conclude $Q(y) \vee Q(z)$ holds, i.e.,
\[
Q(x_1 x_2 \cdots x_n) \ \vee\ Q(x_{n+1}).
\]
Now consider the two cases:
\begin{itemize}
\item If $Q(x_{n+1})$ holds, then in particular $Q(x_1) \vee \cdots \vee Q(x_n) \vee Q(x_{n+1})$ holds.
\item If $Q(x_1 x_2 \cdots x_n)$ holds, then by the inductive hypothesis $P(n)$ applied to $x_1,\dots,x_n$, we get
\[
Q(x_1 x_2 \cdots x_n) \rightarrow Q(x_1) \vee \cdots \vee Q(x_n),
\]
so $Q(x_1) \vee \cdots \vee Q(x_n)$ holds, and hence so does $Q(x_1) \vee \cdots \vee Q(x_n) \vee Q(x_{n+1})$.
\end{itemize}
In either case, $Q(x_1) \vee \cdots \vee Q(x_n) \vee Q(x_{n+1})$ holds, which is exactly $P(n+1)$.

By mathematical induction, $P(n)$ holds for all integers $n \ge 1$.
\end{proof}

\medskip
\noindent\emph{Remark (fixing a gap).} The original argument substituted $y$ and $z$ back into $Q(y) \vee Q(z)$ and asserted the result followed ``by the inductive hypothesis,'' but never explicitly invoked $P(n)$ on $Q(y) = Q(x_1 \cdots x_n)$. The step above makes that invocation explicit: we need $P(n)$ to expand $Q(x_1\cdots x_n)$ into $Q(x_1) \vee \cdots \vee Q(x_n)$.

\bigskip
\partlabel{b}
\emph{Theorem 1.} If the product of $n \ge 1$ real numbers is irrational, then at least one of those numbers must be irrational.

To derive this from part (a), define
\[
Q(x) := \text{``$x$ is irrational.''}
\]
Before part (a) applies to this $Q$, we must check $Q$ satisfies the hypothesis of part (a), namely the base property
\[
\forall x, y \in \R.\ \big( Q(xy) \rightarrow Q(x) \vee Q(y) \big).
\]
Equivalently (by contrapositive), we must prove the following lemma:

\begin{quote}
\textbf{Lemma.} For all $x, y \in \R$, if $x$ and $y$ are both rational, then $xy$ is rational.
\end{quote}

Once the Lemma is established, part (a) (with this $Q$) gives exactly Theorem 1.

\medskip
\partlabel{c}
\emph{Proof of the Lemma, by contradiction.}
\begin{proof}
Suppose, toward a contradiction, that $x$ and $y$ are both rational but $xy$ is irrational. Since $x$ and $y$ are rational, by definition there exist $p_1, q_1, p_2, q_2 \in \Z$ with $q_1 \ne 0$, $q_2 \ne 0$, such that
\[
x = \frac{p_1}{q_1}, \qquad y = \frac{p_2}{q_2}.
\]
Then
\[
xy = \frac{p_1}{q_1}\cdot\frac{p_2}{q_2} = \frac{p_1 p_2}{q_1 q_2}.
\]
Since $p_1, p_2, q_1, q_2 \in \Z$, we have $p_1 p_2 \in \Z$ and $q_1 q_2 \in \Z$; moreover $q_1 q_2 \ne 0$ since $q_1 \ne 0$ and $q_2 \ne 0$. Hence $xy$ is expressible as a ratio of two integers with nonzero denominator, i.e.\ $xy$ is rational. This contradicts our assumption that $xy$ is irrational. Therefore, no such $x, y$ exist: whenever $x$ and $y$ are both rational, $xy$ is rational.
\end{proof}

\emph{Finishing the proof of Theorem 1.} The Lemma is precisely the contrapositive of the required property of $Q(x) := $ ``$x$ is irrational,'' namely
\[
\forall x,y \in \R.\ Q(xy) \rightarrow Q(x) \vee Q(y).
\]
(Indeed: if $xy$ is irrational, then $x,y$ cannot both be rational, by the Lemma; so at least one of $x, y$ is irrational.) Since $Q$ satisfies the hypothesis of part (a), part (a) tells us that for every integer $n \ge 1$ and all $x_1, \dots, x_n \in \R$,
\[
Q(x_1 x_2 \cdots x_n) \rightarrow Q(x_1) \vee \cdots \vee Q(x_n),
\]
i.e., if $x_1 x_2 \cdots x_n$ is irrational, then at least one of $x_1, \dots, x_n$ is irrational. This is exactly Theorem 1. $\blacksquare$

\end{document}
